% Glava dokumenta
\documentclass[slovene, usenames,dvipsnames]{beamer}
\usepackage[utf8x]{inputenc}
\usepackage{lmodern}
\usepackage[T1]{fontenc}

% usetheme is the same as usepackage
% it just automatically prepends the 'beamertheme'
%\usetheme[]{progressbar}
\usefonttheme{progressbar}
\useoutertheme{progressbar}
\useinnertheme{progressbar}
\progressbaroptions{ headline=sections, titlepage=normal, frametitle=normal}
\usecolortheme{seahorse}
%\usecolortheme{progressbar}

% \usetheme[sectionpage=none, subsectionpage=none, numbering=none, progressbar=frametitle]{metropolis} % izbira nastavitev oblike in barv strani
%\usecolortheme{seahorse}
\usepackage{tikz}
\usepackage[slovene]{babel}
%\usetikzlibrary{arrows}
\usepgflibrary{arrows}% for more options on arrows
\usetikzlibrary{arrows.meta}
\usetikzlibrary{positioning}
\usepackage{varwidth}
\usepackage{color}
\usepackage{amsmath}
\usepackage{amsthm}
\usepackage{xcolor}

%\setbeamercolor{frametitle}{
%  use=palette primary,
%  bg=RoyalBlue}

%\setbeamercolor{background canvas}{
% use=palette primary,
% bg=White}

%\setbeamercolor{progress bar}{
%  fg=Blue}

%\setbeamercolor{progress bar in head/foot}{
%  fg=Red}

\usepackage{varwidth}
\setlength{\abovedisplayskip}{0pt}
\setlength{\belowdisplayskip}{0pt}


\uselanguage{slovene}
\languagepath{slovene}
\deftranslation[to=slovene]{Definition}{Definicija}
\usepackage{outlines}
\usepackage[document]{ragged2e}
\usepackage[customcolors,norndcorners]{hf-tikz}
\usetikzlibrary{calc} %% added

\graphicspath{{./slike/}{./slike/fazni_diagrami/}{../eps_pdf/}}
\DeclareGraphicsExtensions{.eps,.jpeg,.png,.gif,.pdf}
\usepackage[outdir=../uporabljene_slike/]{epstopdf}
\epstopdfsetup{
	suffix=,
}

\newcommand{\Alpha}{A}
\newcommand{\Beta}{B}
\newcommand{\Epsilon}{E}
\newcommand{\Kappa}{K}
\newcommand{\dif}{\mathrm{d}}

\setbeamertemplate{navigation symbols}{}
\setbeamertemplate{section in toc}[sections numbered]
\setbeamertemplate{subsection in toc}[subsections numbered]

\setlength{\abovedisplayskip}{0pt}
\setlength{\belowdisplayskip}{0pt}

\author{Jure Lapajne}
\title{Calculation of nmr parameters in paramagnetic metal-organic materials}

\begin{document}
\titlepage

\begin{frame}{Presentation plan}
  \begin{enumerate}
  \item MOFs = Metal---organic frameworks \pause
  \item NMR --- what and why \pause
  \item NMR parameter calculation \pause
  \item DFT \pause
  \item Preliminary results
  \end{enumerate}
\end{frame}

\section{MOFs}
\begin{frame}{MOF = METAL-ORGANIC FRAMEWORK }
  \begin{minipage}[]{0.5\textwidth}
    \begin{itemize}
      \onslide<2-> \item crystal structure: central metallic ions + organic ligands 
      \onslide<3-> \item diversity: large number of possible combinations  
      \onslide<4-> \item wide usage: gas storage, clean energy applications,
      nonlinear optics, catalysts
      \onslide<5-> \item metallic ions commonly feature unpaired electrons
      \onslide<6-> \item NMR spectra feature large shifts caused by unpaired electrons
    \end{itemize}
  \end{minipage}%
\onslide<2->  \begin{minipage}[]{0.5\textwidth}
    \begin{minipage}[]{\textwidth}
      \centering
      \includegraphics[width=0.6\textwidth]{hkust.png}
    \end{minipage}
    
    \begin{minipage}[]{\textwidth}
      \centering
      \begin{tikzpicture}[remember picture,overlay]
      
        \def \dx {-1.2cm}
        \def \dy {-0.2cm}
        \def \Sdim {0.2cm}
        \definecolor{FillColor1}{RGB}{102,178,255};
        \definecolor{FillColor1}{RGB}{102,178,255};
        \def \ygap {0.5cm}
        \def \xgap {1.5cm}
        \def \radius {0.07}

        \draw [fill=black, thick, fill opacity=0.5] ($(\dx, \dy)$) circle [radius=\radius] node[right, black] {carbon};
     
        \draw [fill=red, red, thick] ($(\dx, \dy) + (0, -\ygap)$) circle [radius=\radius] node[right, black, align=center] {oxygen};

        \draw [fill=brown, thick, orange] ($(\dx, \dy) + (\xgap,  0.0)$) circle [radius=\radius] node[right, black] {copper};

        \draw [thick, draw=black, fill=white] ($(\dx, \dy) + (\xgap, -\ygap)$) circle [radius=\radius] node[right, black] {hydrogen};
      \end{tikzpicture}
    \end{minipage}
  \end{minipage}  
\end{frame}

\section{NMR}
\begin{frame}{Nuclei in external magnetic field}
  \begin{itemize}
  \onslide<2-> \item strong external magnetic field: several T
  \onslide<3->\item nuclei with magnetic moment: lowest energy state splits
  \onslide<4->  \item two new states $\Delta E$ apart
  \onslide<5-> \item radio frequency spectrum: excitations from low to high energy states
  \onslide<6-> \item absorption peak at $\Delta E=\hbar \omega_{res}$
    \item $\omega_{res}$ depends on $B_{eff}(observed nucleus)$
    \end{itemize}
    \onslide<3->   \begin{minipage}[]{\textwidth}
      \centering
    \includegraphics[width=0.5\textwidth]{spin_energy_states.png}
\end{minipage}
\end{frame}


\begin{frame}{Effective magnetic field}
  Several parameters affect $B_{eff}(nucleus)$ and $\omega_{res}$: \pause
  \begin{itemize}
  \item electronic structure -- shielding of external magnetic field \pause 
  \item spin--spin coupling to nearby nuclei and unpaired electrons \pause
  \item unpaired electrons: large paramagnetic shifts \pause
    \item calculation of nmr parameters: accurate knowledge of electronic wavefunction
   \end{itemize}
 \begin{minipage}[]{0.4\textwidth}
      \centering
    \includegraphics[width=0.8\textwidth]{cuacac_spin_density.png}
 \end{minipage}%
 \begin{minipage}[]{0.6\textwidth}
   \centering
    \begin{itemize}
    \item MOFs: metal-organic frameworks \pause
    \item transition metal atoms: unpaired electrons \pause
   \end{itemize}
 \end{minipage}
\end{frame}


\begin{frame}{Chemical and hyperfine shifts}

     \onslide<2->    \begin{definition}
 Total shift tensor $\underline{\underline{\sigma}}$ is defined by:
        \vspace*{-0.5\baselineskip}
        \centering
    \begin{equation} \nonumber
      \vec{B}_{eff}=\vec{B}_0\left( \underline{\underline{ \mathrm I }} - \underline{\underline{ \sigma}} \right).
    \end{equation}
  \end{definition}

  \centering Two sources of chemical shifts:
  
  \begin{minipage}[t]{0.5\textwidth}
    \begin{itemize}
    \onslide<3-> \item electron density change caused by applied external magnetic field
    \onslide<3-> \item depends on electron density $n(\vec r)$
      \end{itemize}
  \end{minipage}%
  \begin{minipage}[t]{0.5\textwidth}
    \begin{itemize}
    \onslide<3->  \item coupling between unpaired electron and observed nuclei
    \onslide<3-> \item depends on spin density $n_{\uparrow}(\vec r_{nuclei}) - n_{\downarrow}(\vec r_{nuclei})$
      \end{itemize}
    \end{minipage}
    Calculation of NMR parameters:\\\bf{accurate electronic wave function needed!}
  \end{frame}

  


\end{document}